% Suggested computational-appendix insertion. Uses the article's existing macros.
% All counts below are populated only after the complete CSV coverage audit.
% Add inside the existing bibliography:
% \bibitem{MMP2025} M. Mula, N. Murru and F. Pintore, On random sampling of
% supersingular elliptic curves, Ann. Mat. Pura Appl. 204 (2025), 1293--1335.
% \href{https://doi.org/10.1007/s10231-024-01528-x}{doi:10.1007/s10231-024-01528-x}.
\subsection{Direct verification at square levels}
\label{app:direct-square-levels}

The standalone PARI/GP program \texttt{verify\_square\_direct.gp} checks
\[
\lambda_p(X_0(\ell^2)^+)=\frac{\ell+1}{2}L(p),\qquad \ell=3,5,7,
\]
by enumerating the supersingular coarse points in characteristic $p$.
It uses no genus formula, class number, modular-form trace or
characteristic-zero fixed-point count. The value $L(p)$ is obtained directly
from the supersingular $j$-values, independently of Deuring's formula.

The supersingular set is generated from the roots in $\mathbb F_{p^2}$ of the
Legendre Hasse polynomial
\[
H_p(z)=\sum_{i=0}^{(p-1)/2}\binom{(p-1)/2}{i}^{\!2}z^i,
\qquad j(z)=256\frac{(z^2-z+1)^3}{z^2(z-1)^2}.
\]
This uses the classical Hasse criterion; see \cite[Proposition 4.19]{MMP2025}.
For each resulting $j$, the program constructs a short Weierstrass model whose
$p^2$-Frobenius is scalar, and enumerates its $\ell+1$ order-$\ell$ subgroup
kernels from its division polynomial. V\'elu's formulas construct the isogenies.
The unique second step whose composite kernel is $E[\ell]$ is discarded.
The remaining $\ell(\ell+1)$ composite kernels are checked to be distinct.
They are exactly the cyclic subgroups of order $\ell^2$.

The full subgroup kernel is retained, rather than only its triple of
$j$-values. Quotienting these kernels by the source automorphisms produces the
geometric coarse points of $X_0(\ell^2)$. In particular the extra automorphisms
at $j=0,1728$ are included. The induced Frobenius $F$ raises coefficients to
their $p$-th powers. Fricke reversal $W$ is constructed from the actual dual
isogenies, with any intermediate canonical-model automorphism corrected.
For every edge the program verifies the complete rational $x$-map identity
$x(\widehat\phi(\phi(P)))=x([\ell]P)$. The sign on the ordinate is irrelevant
to the subgroup kernel. It also checks on every coarse point that
\[
F^2=W^2=1,\qquad FW=WF.
\]
An orbit $\{P,W(P)\}$ is rational precisely when $F(P)$ belongs to that orbit.
The enumeration therefore counts points with Frobenius-swapped lifts as well
as points with rational lifts. A separate finite-set check verifies that the
count is $(\#\Fix(F)+\#\Fix(WF))/2$.

Run under PARI/GP 2.15.4, the enumeration agrees with the formula for every
prime $5\leq p\leq1000$ with $p\ne\ell$, and separately for
$p=1009,2003,5003$, at each of the three levels. There are $505$ distinct
admissible pairs ($169$ at level $9$, and $168$ each at levels $25$ and $49$),
with no mismatches and no skipped admissible cases. The computation constructs
$695{,}714$ cyclic subgroup kernels and $685{,}906$ coarse covering-curve points.
The initial $178$ comparisons through $300$ also agree on rerunning.
A second enumeration uses the middle curve and an unordered pair of distinct
order-$\ell$ subgroups; it shares the initial Hasse-set and prime-order kernel
constructors, but uses no V\'elu or dual-map reversal routines. It agrees in all
$505$ cases and checks $5{,}260$ rational middle-invariant fibres.
The CSV, kernel and permutation snapshots, source code and coverage audit
are supplied with the verification package.

At $p=13$ the only supersingular invariant is $j=5$, so $L(13)=1$.
The complete direct enumeration gives
\begin{center}
\begin{tabular}{rrrrr}\toprule
$\ell^2$ & $X_0$ points & Fricke orbits & Rational orbits & Swapped lifts\\\midrule
9 & 12 & 6 & 2 & 2\\
25 & 30 & 15 & 3 & 3\\
49 & 56 & 28 & 4 & 3\\\bottomrule
\end{tabular}
\end{center}
Every path here has the same triple $(5,5,5)$ of $j$-values, which demonstrates
why collapsing paths to distinct triples would give an incorrect count.
The package includes every kernel and both permutations for these examples,
and a separate audit that reloads and recounts them.

These computations test the rational supersingular point count directly.
They do not compute a normalized Hauptmodul polynomial and do not replace the
geometric proof or address additional moonshine normalizer actions.
